\FloatBarrier
\appendix

\section{Types: Session-Type Laws and Well Formedness Predicates}
\label{sec:appendix-types}

\begin{figure}[htbp]
  \begin{mathpar}
    \RuleUUnit \and
    \RuleUArr \and
    \RuleUProd \and
    \RuleUVariant \\
    \RuleUCtxEmpty \and
    \RuleUCtxVar \and
    \RuleUCtxSeq \and
    \RuleUCtxPar \\
    \RuleSNewSkip \and
    \RuleSNewSend \and
    \RuleSNewRecv \and
    \RuleSNewVar \\
    \RuleSNewSeq \and
    \RuleSNewBranch \and
    \RuleSNewChoice \and
    \RuleSNewRec
  \end{mathpar}
  \caption{Unrestricted types and contexts ($\Unr\T$, $\Unr\Ctx$); session
    types valid for $\CNew$ ($\SNew\S$)}
  \label{fig:omitted-predicates}
\end{figure}

\begin{figure}[htbp]
  \medskip
  Laws for context-free session types.
  \begin{mathpar}
    \SSeq\SSkip\S = \S \and
    \SSeq \S\SSkip = \S \and
    \SSeq{\S[1]}{(\SSeq{\S[2]}{\S[3]})} =
    \SSeq{(\SSeq{\S[1]}{\S[2]})}{{\S[3]}} \and
    \SRec\SVar\S = \S{[\SRec\SVar\S/\SVar]} \and
    \SSeq{(\SChoice{\ell:\S[\ell]}{\ell\in L})}\S =
    \SChoice{\ell:\SSeq{\S[\ell]}\S}{\ell\in L} \and
    \SSeq{(\SBranch{\ell:\S[\ell]}{\ell\in L})}\S = \SBranch{\ell:\SSeq{\S[\ell]}\S}{\ell\in L}
  \end{mathpar}

  \medskip
  Laws for duality $\SDual\S$ (there is no dual for $\SDrop$ and $\SAcq$).
  \begin{mathpar}
    \SDual\SSkip = \SSkip \and
    \SDual{(\SSeq{\S[1]}{\S[2]})} = \SSeq{\SDual{\S[1]}}{\SDual{\S[2]}}
    \and
    \SDual{\SSend\T} = \SRecv\T \and
    \SDual{\SRecv\T} = \SSend\T \and
    \SDual{(\SChoice{\ell:\S[\ell]}{\ell\in L})} =
    \SBranch{\ell:\SDual{\S[\ell]}}{\ell\in L} \and
    \SDual{(\SBranch{\ell:\S[\ell]}{\ell\in L})} =
    \SChoice{\ell:\SDual{\S[\ell]}}{\ell\in L} \and
    \SDual\STerm = \SWait \and
    \SDual\SWait = \STerm \and
    \SDual{(\SRec\SVar\S)} = \SDual{\S{[\SRec\SVar\S/\SVar]}}
  \end{mathpar}
  \caption{Session-type equality and duality}
  \label{fig:session-laws}
\end{figure}

\FloatBarrier

\section{Dynamics: Congruence and Reduction Rules}

\begin{figure}[tp]
\begin{align*}
  \U,\V \grmeq
    & \C
      \grmor \EVar
      \grmor \EAbs \EVar\E
      \grmor \EInj{i} \V
      \grmor \EPair \U\V
      % \grmor \CSend \V 
              \tag{Values}
  \\
  \EC \grmeq&
              \ECEmpty
              \grmor \EAppD[\DirLeft] \EC\E
              \grmor \EAppD[\DirLeft] \V\EC
              \grmor \EAppD[\DirRight] \E\EC
              \grmor \EAppD[\DirRight] \EC\V
              \grmor \ESeq \EC\E
              \grmor \EPair \EC\E
              \grmor \EPair \V\EC
              \\ \grmor
    & \EInj{i}\EC \grmor \ELet {\EPair \EVar\EVar}\EC\E
              \grmor \ECase \EC {\ECaseBranch* i \EVar \E}
              \tag{Evaluation contexts}
  \\
  \FC \grmeq
  & \PExp \EC
    \tag{Thread evaluation contexts}
  \\
  \CC \grmeq
  & \ECEmpty
    \grmor \PPar{\CC}{\Proc}
    \grmor \PNew \BindGroup\BindGroup \CC
    \tag{Process contexts}
\end{align*}
  \caption{Values and evaluation contexts}
  \label{fig:values-evaluation-contexts}
\end{figure}

\begin{figure}[htbp]
  \begin{mathpar}
    \RuleExpRedApp \and
    \RuleExpRedSeq \and
    \RuleExpRedPairElim \and
    \RuleExpRedUnfold \and
    \RuleExpRedVariantElim \and
    \RuleExpRedCtx
  \end{mathpar}
  \caption{Expression reduction ($\E \EReduce \E$)}
  \label{fig:expression-reduction}
\end{figure}

\begin{figure}[htbp]
  \begin{mathpar}
    % \RuleProcCongRefl \and
    % \RuleProcCongTrans \and
    \RuleProcCongParComm \and
    \RuleProcCongParAssoc \and
    \RuleProcCongParUnit \and
    \RuleProcCongResSwap \and
    \RuleProcCongResComm \and
    \RuleProcCongExtend
  \end{mathpar}
  \caption{Process congruence}
  \label{fig:configuration-congruence}
\end{figure}

\FloatBarrier

\section{Translation: Congruence of Translated Processes}


\begin{figure}[htbp]
  \begin{mathpar}
    %\RuleTransCongRefl \and
    \RuleTransCongParComm \and
    \RuleTransCongParAssoc \and
    \RuleTransCongParUnit \and
    \RuleTransCongResSwap \and
    \RuleTransCongResComm \and
    \RuleTransCongFlagComm \and
    \RuleTransCongResFlagComm \and
    \RuleTransCongResExtend \and
    \RuleTransCongFlagExtend
    %\RuleTransCongTrans
  \end{mathpar}
  \caption{Congruence rules of translated processes}
  \label{fig:translated-congruence}
\end{figure}

\FloatBarrier

\section{Algorithmic Typing}
\label{sec:algorithmic-typing-appendix}

\begin{figure}[htbp]
  \begin{minipage}{0.49\linewidth}
  \[
    \Ctx[1] \CtxJoin \Ctx[2] = \begin{cases}
      \Ctx[1] \CtxPar \Ctx[2] & \text{if $\Dir \in \{ \DirUnord, \DirNone \}$} \\
      \Ctx[1] \CtxSeq \Ctx[2] & \text{if $\Dir = \DirLeft$} \\
      \Ctx[2] \CtxSeq \Ctx[1] & \text{if $\Dir = \DirRight$}
    \end{cases}
  \]
  \end{minipage}
  \begin{minipage}{0.49\linewidth}
  \[
    \CtxRestrict*{X} = \begin{cases}
      \CtxRestrict*[1]{X} \CtxPar \CtxRestrict*[2]{X} & \text{if $\Ctx=\Ctx[1] \CtxPar \Ctx[2]$} \\
      \CtxRestrict*[1]{X} \CtxSeq \CtxRestrict*[2]{X} & \text{if $\Ctx=\Ctx[1] \CtxSeq \Ctx[2]$} \\
      \CtxVar \EVar \T & \text{if $\Ctx=\CtxVar \EVar T$ and $\EVar \in X$} \\
      \CtxEmpty & \text{otherwise}
    \end{cases}
  \]
  \end{minipage}
  \[
    \CaseJoinDir(\Ctx, X) = \begin{cases}
      \DirUnord & \text{if $\CtxRestrict*{X} \CtxPar \CtxRestrict*{X^C} \SubCtx \Ctx$}\\
      \DirLeft & \text{if $\CtxRestrict*{X} \CtxSeq \CtxRestrict*{X^C} \SubCtx \Ctx$}\\
      \text{undef.} & \text{otherwise}\\
    \end{cases}
  \]
  \caption{Context join $\Ctx[1] \CtxJoin \Ctx[2]$ joins two contexts
    according to the direction $\Dir$, context restriction
    $\CtxRestrict*{X}$ restricts a context to a set of variables $X$, and $\CaseJoinDir$ calculates if a context $\Ctx$ restricted to a set $X$ can be joined parallelly with $\Ctx$ restricted to the complement set.}
  \label{fig:context-operators}
\end{figure}

\begin{figure}[htbp]
  \begin{mathpar}
    \RuleCSPair \and
    \RuleCSSum \and
    \RuleCSUnit \and
    \RuleCSArrow \and
    \RuleCSContext
  \end{mathpar}
  \caption{Constraint simplification}
  \label{fig:constraint-simplification}
\end{figure}

\subsection{Type Normalisation}
\begin{figure}
  \small
  \hfill \fbox{$\ConcatType {\T[in]} {\T[in]} = \T[out]$}
  \begin{mathpar}
    \TypeConcatenationSkip

    \TypeConcatenationSemiDist

    \TypeConcatenationNonSemi
  \end{mathpar}

  \hfill \fbox{$\NormaliseType {\T[in]} {\T[out]}$}
  \begin{mathpar}
    \TypeNormalisationSemiSkips
    
    \TypeNormalisationSemi

    \TypeNormalisationMu

    \TypeNormalisationChoiceSemi

    \TypeNormalisationBranchSemi

    \TypeNormalisationIdentity
  \end{mathpar}
  \caption{Type Normalisation}
  \label{fig:typing-normalisation}
\end{figure}


\subsection{Implementation}
\label{sec:appendix-implementation}
\subsubsection{Constraint Checking with FreeST}

In the implementation, we use FreeST \cite{DBLP:journals/corr/abs-1904-01284,DBLP:journals/iandc/AlmeidaMTV22, almeida26:_frees_concur_progr_languag_with} version 5.0 to perform equivalence checking. To achieve this, we generate little FreeST programs that typechecks iff the equivalence constraints hold.

For every equivalence constraint $\T[1] \CEq \T[2]$, we generate the following FreeST program:

\begin{lstlisting}[numbers=left]
type T1 : 1S
type T1 = ...

type T2 : 1S
type T2 = ...

id : T1 -> T2
id x = x
\end{lstlisting}

FreeST 5.0 requires specifying kinds for type definitions. 
Kind \lstinline|1S| says that the type needs to be a session type. 
Let \lstinline|T1| and \lstinline|T2| are the translated types from the equivalence constraint, 
it is easy to see that this program only succefully type checks, if \lstinline|T1| and \lstinline|T2| are equivalent types.

The main parts of the translation are:
\begin{itemize}
  \item Translating inline $\mu$ recursion to explicit recursion.
  \item Translating mobility and direction annotations on function and pair types.
\end{itemize}

To translate from $\mu$ to explicit recursion, we generate a fresh name and associate the name with the recursion unrolled and translation applied. 
We then return the associated name, instead of the type. +
Thus the translation function $\tofreest{\S}$ (\Cref{fig:freest-translation} ) takes a session type $S$ and returns a FreeST Type $F$ together with list of explicit type definitions.

\begin{figure}[t!]
  \begin{align*}
    \tofreest{\TypeKeyword{Unit}}
      & = [(), \emptyset] \\
    \tofreest{\SSkip}
      & = [\SSkip, []] \\
    \tofreest{\STerm}
      & = [\STerm, []] \\
    \tofreest{\SWait}
      & = [\SWait, []] \\
    \tofreest{\TArr{\Mob}{\Dir, {\Eff}}{\T[1]}{\T[2]}}
      & = (\tofreest{\T[1]} \to \tofreest{\T[2]}) \TPair \SBranch {\ell : \SSkip} {\ell \in \text{dir} \cup \text{mob} \cup \text{eff}} \quad \text{where} \\
      & \text{dir} = \begin{cases}
        \{ \text{left} \} & \text{if } \Dir = \DirLeft \\
        \{ \text{right} \} & \text{if } \Dir = \DirRight \\
        \{ \text{lin} \} & \text{if } \Dir = \DirUnord \\
        \emptyset & \text{if } \Dir = \DirNone \\
      \end{cases} \\
      & \text{mob} = \begin{cases}
        \{ \text{mobile} \} & \text{if } \Mob = \MobMobile \\
        \{ \text{static} \} & \text{if } \Mob = \MobStationary \\
      \end{cases} \\
      & \text{eff} = \begin{cases}
        \{ \text{impure} \} & \text{if } \Eff = \EffImpure \\
        \{ \text{} \} & \text{if } \Eff = \EffImpure \\
      \end{cases} \\
    \tofreest{\SSend\T}
      & = [\SSend{F}, U] \quad \text{where} \\
      & \quad [F, U] = \tofreest{\T} \\
    \tofreest{\SRecv\T}
      & = [\SRecv{F}, U] \quad \text{where} \\
      & \quad [F, U] = \tofreest{\T} \\
    \tofreest{\S[1] ; \S[2]}
      & = [F_1 ; F_2,\, U_1 \cup U_2] \quad \text{where} \\
      & \quad [F_1, U_1] = \tofreest{\S[1]} \\
      & \quad [F_2, U_2] = \tofreest{\S[2]} \\
    \tofreest{\SChoice{\ell:\S[\ell]}{\ell\in L}}
      & = [+\{\ell : \tofreest{F_\ell}\}_{\ell\in L},\, \bigcup_{\ell\in L} U_\ell] \quad \text{where} \\
      & \quad [F_\ell, U_\ell] = \tofreest{\S[\ell]} \quad \text{for every } \ell \in L \\
    \tofreest{\SBranch{\ell:\S[\ell]}{\ell\in L}}
      & = [\&\{\ell : \tofreest{F_\ell}\}_{\ell\in L},\, \bigcup_{\ell\in L} U_\ell] \quad \text{where} \\
      & \quad [F_\ell, U_\ell] = \tofreest{\S[\ell]} \quad \text{for every } \ell \in L \\
    \tofreest{\SDrop}
      & = [\SSend{\TypeKeyword{Drop}},\, \emptyset] \\
    \tofreest{\SAcq}
      & = [\SRecv{\TypeKeyword{Drop}}, \emptyset] \\
    \tofreest{\SRec\SVar\S}
      & = [\SVar',\, U \cup \{\SVar' \mapsto U\}] \quad \text{where } \SVar' \text{ is fresh and} \\
      & \quad [F, U] = \tofreest{\S\,[\SVar'/\SVar]} \\
    \tofreest{\SVar}
      & = [\SVar, \emptyset]
  \end{align*}
  \caption{Translation of session types to FreeST types.}
  \label{fig:freest-translation}
\end{figure}

\subsubsection{Soundness and the Completeness of the translation}

\begin{definition}[Equivalence relation for \thecalculus types]
  The equivalence relation $\simeq$ on \thecalculus types is the least congruence closed under the equivalence rules from \cite{DBLP:conf/icfp/ThiemannV16} extended with rules in \Cref{fig:cfstb-equivalence}.
\end{definition}

\begin{figure}[t!]
  \begin{mathpar}
    \inferrule{
      \T[1] \simeq \T*[1] \\
      \T[2] \simeq \T*[2]
    }{
      \TArr{\Mob}{\Dir,\Eff}{\T[1]}{\T[2]}
      \simeq
      \TArr{\Mob}{\Dir,\Eff}{\T*[1]}{\T*[2]}
    }
    \and
    \inferrule{
      \T[1] \simeq \T*[1] \\
      \T[2] \simeq \T*[2]
    }{
      \T[1] \TPair \T[2]
      \simeq
      \T*[1] \TPair \T*[2]
    }
    \and
    \inferrule{
      \T[1] \simeq \T*[1] \\
      \T[2] \simeq \T*[2]
    }{
      \TSum{\T[1]}{\T[2]}
      \simeq
      \TSum{\T*[1]}{\T*[2]}
    }
  \end{mathpar}
  \caption{Modified rules for $\simeq$}
  \label{fig:cfstb-equivalence}
\end{figure}

\begin{definition}[Substituting definitions]
  Let $F$ be a FreeST session type and $U$ be set of mapping from variables to FreeST session types such that $\text{freevars}(F) \subseteq \text{dom}(U)$. The subsitution of $U$ to $F$, denoted by $F[U]$, is the type where every variable in $U$ that has a definition in $F$ is substituted. A well formed subsitution has no free variables.
\end{definition}

\begin{lemma}[Soundness of the translation]
  If $\T[1] \simeq \T[2]$, then $F_1 [U_1] \simeq F_2 [U_2]$ where $(F_1, U_1) = \tofreest{\T[1]}$ and $(F_2, U_2) = \tofreest{\T[2]}$.
  \label{lem:sound-translation}
\end{lemma}
\begin{proof}
  We sketch a proof. The soundness proof requires two things: the translation of annotations in the arrow and product types preserves equivalence and the equivalence of the explicit recursion with the $\mu$ binder syntax. For the former we can perform induction on the equivalence of $\T[1] \simeq \T[2]$ and use equivalence and bisimulation lemmas from \cite{DBLP:conf/icfp/ThiemannV16} for regular type and session type cases, respectively. The pairs we introduce for annotations can be handled trivially by the type equivalence rules. To show explicit recursion translation preserves equivalence, we can use the laws for $\mu$ types from \cite{DBLP:conf/icfp/ThiemannV16}.
\end{proof}

\begin{lemma}[Completeness of the translation]
  Let $(F_1, U_1) = \tofreest{\T[1]}$ and $(F_2, U_2) = \tofreest{\T[2]}$. If $F_1 [U_1] \simeq F_2 [U_2]$, then $\T[1] \simeq \T[2]$.
\end{lemma}
\begin{proof}
  Similar to \ref{lem:sound-translation} but the opposite direction.
\end{proof}

\FloatBarrier

\section{Metatheory}
\label{sec:appendix-metatheory}

\begin{figure}[htbp]
  \begin{mathpar}
    \RuleExpBlocked \and
    \RuleParBlocked \and
    \RuleNuBlocked
  \end{mathpar}
  where
  \[
    \begin{array}{ll}
      \BlockedChannelsSend & \BlockedChannelsRecv \\
      \BlockedChannelsSelect & \BlockedChannelsBranch \\
      \BlockedChannelsLSplit & \BlockedChannelsRSplit \\
      \BlockedChannelsFork
    \end{array}
  \]
  \caption{Blocked Expressions (Blocked $\E$) and Blocked Processes (Blocked $\P$)}
  \label{fig:blocked-expr-proc}
\end{figure}

\subsection{Type Safety for Declarative Typing}
\label{sec:appendix-metatheory-ts}

\providecommand{\proofcase}[1]{\smallskip\noindent\textit{Case}~#1\quad}

We give the proof of process progress
(\cref{thm:process-progress}). The argument is the concurrent analogue of
the sequential progress proof for expressions: an inductive walk over the
process structure that, at every restriction, either exhibits a reduction or
certifies that the configuration is genuinely deadlocked. The blocked-process
predicate $\Blocked\Proc$ and the active-channel function $\BC(\Proc)$ used
below are the ones from \cref{fig:blocked-expr-proc}; the evaluation and
process contexts are those of \cref{fig:values-evaluation-contexts}.

\paragraph{Preliminaries.}
Throughout, $\SessCtx$ denotes a \emph{session context}: a context all of whose
entries are session (channel) types. In such a context every variable denotes a
channel endpoint. This is the intrinsic side condition on the value predicate:
a variable counts as a value only under a channel context, so canonical-form
reasoning for values is available exactly when the context is a session
context.

\begin{definition}[Poised thread, active constant]
  \label{def:poised}
  A thread $\PExp\E$ is \emph{poised} when $\E = \ECapp{\EApp\C\V}$ for some
  constant $\C$ and value $\V$, i.e.\ when $\Blocked\E$ holds by
  \textsc{B-ExpBlocked}. We call $\EApp\C\V$ its \emph{active redex}, $\C$ its
  \emph{active constant}, and, when $\V$ (or its relevant component) is a
  channel variable $\EVar$, we say the thread is \emph{poised on} $\EVar$.
\end{definition}

We partition the constants by the shape of the reduction they participate in.
\begin{itemize}
  \item \emph{Spawning} constants $\CNew,\CFork$: their active redex reduces at
    the thread level, independently of any restriction, by \textsc{R-New}
    resp.\ \textsc{R-Fork}. They contribute $\BC = \emptyset$.
  \item \emph{Communicating} constants
    $\CSend,\CRecv,\CSelect{},\CBranch,\CTerm,\CWait$: their active redex fires
    only against a matching action on the \emph{dual} endpoint, by
    \textsc{R-Com}, \textsc{R-Choice}, resp.\ \textsc{R-Close}. Each contributes
    $\BC = \{\EVar\}$ for its channel operand $\EVar$.
  \item \emph{Asynchronous} constants $\CDrop,\CDiscard,\CSplit,\CSSplit$: their
    active redex fires alone once the operand channel is bound by an enclosing
    restriction, by \textsc{R-Drop}, \textsc{R-Discard}, \textsc{R-LSplit},
    resp.\ \textsc{R-RSplit}; no partner thread is required. They contribute
    $\BC = \emptyset$.
  \item \emph{Acquire} $\CAcq$ is the hybrid: it synchronizes with an
    \emph{earlier} $\CDrop$ on the dual side through the one-bit flag recorded in
    the binder. Its redex fires by \textsc{R-Acquire} exactly when that flag is
    set --- syntactically, when the operand channel carries the leading
    separator marker $\BLeadSep{}$ in the bind group --- and contributes
    $\BC = \{\EVar\}$ otherwise.
\end{itemize}

\begin{lemma}[Thread progress]
  \label{lem:thread-progress}
  If $\HasType[\EffImpure]\SessCtx\E\BUnit$, then either $\E$ is a value,
  or $\E \EReduce \E*$ for some $\E*$, or $\E$ is poised. Moreover, if $\E$ is a
  value it is $\CUnit$.
\end{lemma}
\begin{proof}
  The three-way alternative is expression progress
  (\cref{thm:expression-progress}), noting that a stuck redex $\ECapp{\EApp\C\V}$
  is exactly a poised thread. For the last clause: a value of type $\BUnit$ in a
  session context is, by canonical forms, the constant $\CUnit$ (the only value
  introduction for $\BUnit$; a variable would have a session type, not
  $\BUnit$).
\end{proof}

\begin{lemma}[Synchronization]
  \label{lem:sync}
  Let $\PNew{\BSeq\EVar{\BindGroup[1]}}{\BSeq\EVary{\BindGroup[2]}}\Proc$ bind the
  dual endpoints $\EVar : \SSeq\S{\S*}$ and $\EVary : \SDual{(\SSeq\S{\S*})}$,
  and suppose one thread of $\Proc$ is poised on $\EVar$ with a communicating
  active constant and another is poised on $\EVary$. Then the two active
  constants are dual --- one of the pairs $(\CSend,\CRecv)$,
  $(\CSelect{k},\CBranch)$, $(\CTerm,\CWait)$ --- and the restriction reduces by
  \textsc{R-Com}, \textsc{R-Choice}, resp.\ \textsc{R-Close}.
\end{lemma}
\begin{proof}
  The constant-typing rules force the operand session type of each poised
  redex to be the head of the endpoint's session type. Since the two endpoints
  are dual and duality maps $\SSend\T \mapsto \SRecv\T$,
  $\SChoice{\ell:\S[\ell]}{} \mapsto \SBranch{\ell:\SDual{\S[\ell]}}{}$ and
  $\STerm \mapsto \SWait$ (\cref{fig:session-laws}), the two heads are dual, so
  the active constants match one synchronization rule, whose left-hand side is
  precisely the given pair of poised redexes.
\end{proof}

\begin{theorem}[Process progress, restated]
  \label{thm:process-progress-app}
  If $\ProcHasType\SessCtx\Proc$, then either $\Proc \Congruent \PExp\CUnit$, or
  $\Blocked\Proc$, or there exists $\Proc*$ with $\Proc \PReduce \Proc*$.
\end{theorem}
\begin{proof}
  By induction on the derivation of $\ProcHasType\SessCtx\Proc$. The subsumption
  rule \textsc{TP-Weaken} (which only relaxes the context ordering) is discharged
  by the induction hypothesis, since neither $\PReduce$, $\Congruent$, nor
  $\Blocked{-}$ depends on the ordering slack. We treat the three structural
  rules; recall that $\Blocked{-}$ and $\Congruent$ are read up to the pruning of
  terminated $\PExp\CUnit$ threads via \textsc{Par-Unit}, which we use silently
  below.

  \proofcase{\textsc{TP-Expr}, $\Proc = \PExp\E$ with $\HasType[\EffImpure]\SessCtx\E\BUnit$.}
  Apply \cref{lem:thread-progress}.
  \begin{itemize}
    \item $\E$ is a value: then $\E = \CUnit$, so $\Proc = \PExp\CUnit$; the first
      alternative holds.
    \item $\E \EReduce \E*$: then $\Proc \PReduce \PExp{\E*}$ by \textsc{R-Exp};
      the third alternative holds.
    \item $\E = \ECapp{\EApp\C\V}$ is poised. If $\C \in \{\CNew,\CFork\}$, the
      thread reduces by \textsc{R-New} resp.\ \textsc{R-Fork} (third
      alternative). Otherwise $\C$ acts on a channel operand $\EVar$ which, as
      $\Proc$ has no restriction, is free, hence bound in $\SessCtx$. No process
      rule applies to a channel that no restriction binds --- every channel rule
      \textsc{R-Com}, \dots, \textsc{R-Acquire} discharges a binder. Thus
      $\Blocked\Proc$ by \textsc{B-ExpBlocked} (second alternative).
  \end{itemize}

  \proofcase{\textsc{TP-Par}, $\Proc = \PPar{\Proc[1]}{\Proc[2]}$.}
  Apply the induction hypothesis to $\Proc[1]$ and $\Proc[2]$.
  \begin{itemize}
    \item If $\Proc[i] \PReduce \Proc[i]'$ for some $i$, then $\Proc \PReduce$ by
      \textsc{R-Par} (for $i=2$, precede it with $\PPar{\Proc[1]}{\Proc[2]}
      \Congruent \PPar{\Proc[2]}{\Proc[1]}$ and close under \textsc{R-Struct}).
    \item Otherwise neither component reduces. If $\Proc[i] \Congruent
      \PExp\CUnit$, then $\Proc \Congruent \Proc[3-i]$ by \textsc{Par-Unit}, and
      we conclude from the alternative obtained for $\Proc[3-i]$.
    \item Otherwise both components are blocked (each is, up to pruning, a
      non-empty parallel composition of poised threads and blocked
      restrictions), so $\Blocked\Proc$ by \textsc{B-ParBlocked}.
  \end{itemize}

  \proofcase{\textsc{TP-Res}, $\Proc = \PNew{\BindGroup[1]}{\BindGroup[2]}\ProcQ$.}
  Inversion of \textsc{TP-Res} types $\ProcQ$ in $\SessCtx$ extended by the two
  bind groups, which introduce the freshly bound dual endpoints; the extended
  context is again a session context, so the induction hypothesis applies to
  $\ProcQ$.
  \begin{itemize}
    \item $\ProcQ \PReduce \ProcQ'$: then $\Proc \PReduce
      \PNew{\BindGroup[1]}{\BindGroup[2]}{\ProcQ'}$ by \textsc{R-Bind}.
    \item $\ProcQ \Congruent \PExp\CUnit$: this cannot arise while
      $\BindGroup[1],\BindGroup[2]$ bind unconsumed endpoints. Their session
      types $\SSeq\S\STerm$ and $\SDual{(\SSeq\S\SWait)}$ are linear, and
      $\BindGivesCtx{}{}{}$ forces every endpoint to be consumed by exactly one
      thread; a body congruent to a single $\CUnit$ thread leaves them unused
      and is untypable. (The terminal $\CUnit$ configuration is produced only
      \emph{by} \textsc{R-Close}, which discharges the restriction in the same
      step, so it is never the body \emph{under} a live restriction.)
    \item $\ProcQ$ is blocked. This is the crux. Up to pruning, $\ProcQ$ is a
      parallel composition of poised threads and blocked inner restrictions; the
      inner restrictions are already handled by the induction hypothesis, so the
      only channels not yet accounted for are the ones bound \emph{here}. We
      analyze the poised threads that act on them.
      \begin{enumerate}
        \item If some thread is poised with an asynchronous constant
          ($\CDrop,\CDiscard,\CSplit,\CSSplit$) on a channel bound here, or with
          $\CAcq$ on a bound channel whose flag is set (leading separator
          $\BLeadSep{}$ present), then --- after rotating that channel to the head
          of the bind group by \textsc{Res-Swap}/\textsc{Res-Comm} under
          \textsc{R-Struct} --- the restriction reduces by \textsc{R-Drop},
          \textsc{R-Discard}, \textsc{R-LSplit}, \textsc{R-RSplit}, resp.\
          \textsc{R-Acquire}. (These are exactly the constants with $\BC =
          \emptyset$, whose steps need no partner: an endpoint poised on one of
          them never blocks its restriction.)
        \item Otherwise, if some channel bound here has its two dual endpoints
          simultaneously poised at communicating constants, then by
          \cref{lem:sync} the restriction reduces by \textsc{R-Com},
          \textsc{R-Choice}, resp.\ \textsc{R-Close} (again rotating that channel
          to the head first, if needed).
        \item Otherwise, every redex poised on a channel bound here is either a
          communicating action whose dual endpoint is \emph{not} poised (its
          partner is a value thread or is itself blocked on a different channel),
          or an $\CAcq$ whose flag is unset. In both cases the head pair
          $(\EVar,\EVary)$ of the binder satisfies $\EVar \notin \BC(\ProcQ)$ or
          $\EVary \notin \BC(\ProcQ)$; since any bound channel can be rotated to
          the head by \textsc{Res-Swap}/\textsc{Res-Comm} while preserving both
          $\Blocked{-}$ and the failure of synchronization, the premise of
          \textsc{B-NuBlocked} holds. With $\Blocked\ProcQ$ from the induction
          hypothesis, \textsc{B-NuBlocked} gives $\Blocked\Proc$.
      \end{enumerate}
  \end{itemize}
  This exhausts the cases and closes the induction.
\end{proof}

\noindent\emph{Remark (role of the session-context hypothesis).}
The hypothesis that $\SessCtx$ is a session context is used twice, both in the
\textsc{TP-Expr} case: to conclude $\E = \CUnit$ from ``$\E$ is a value of type
$\BUnit$'', and to guarantee that the operand of a stuck communicating redex is a
\emph{channel} (so that the stuck thread is a genuine communication block rather
than an open term of some other type). Dropping the hypothesis would admit free
variables of arbitrary type and break both steps.

\subsection{Soundness and Completeness for Algorithmic Typing}
\label{sec:appendix-metatheory-sc}

\FloatBarrier

